\documentclass[10pt,a4paper]{amsart}
\usepackage[margin=19mm]{geometry}
\usepackage{amsmath,amssymb,mathtools}
\usepackage[scr=rsfs,cal=euler]{mathalfa}
\usepackage{xfrac}
\usepackage[unicode,hidelinks,bookmarks=false]{hyperref}
\newcommand{\ie}{i.e.\ }
\newcommand{\cf}{cf.\ }
\newcommand{\resp}{respectively\ }
\newcommand{\Subsection}{\subsection}
\providecommand{\nobreakdash}{\nobreak\hbox{-}}
\setlength{\emergencystretch}{2em}
\multlinegap0pt
\usepackage{bm}
\usepackage{bbm}
\usepackage{dsfont}
\usepackage{xspace}
\usepackage{cancel}
\usepackage{mathtools}
\usepackage{amsthm}
\makeatletter
\@ifundefined{prop}{\newtheorem{prop}{Prop}}{}
\makeatother
\title[Single conjugacy classes of isometries in orthogonal groups ]{Single conjugacy classes of isometries in orthogonal groups over local fields}
\author{Fei Xu, Bo Zhang}
\date{Source: arXiv:2602.20481v1 (CC BY 4.0)}
\begin{document}
\maketitle

\noindent\emph{Source and licence.} Single conjugacy classes of isometries in orthogonal groups over local fields. Version arXiv:2602.20481v1 (CC BY 4.0), distributed under the Creative Commons Attribution 4.0 International licence, as recorded in the arXiv licence metadata for this exact version and shown on the abstract page https://arxiv.org/abs/2602.20481v1. Full rights notice: licenses/arxiv-2602.20481-CC-BY-4.0.txt.\par\smallskip

\noindent\textit{Editorial scope.} Counted unit: the Prop whose source label is \texttt{orth-base} in the article (source0.tex lines 174--177, source label \texttt{orth-base}), delivered together with the article's own complete original proof of it (source0.tex lines 178--207, one \texttt{proof} environment of 30 lines). No mathematics is written, repaired or re-derived: statement and proof are the article's own text line for line. Required context retained uncounted: non-degenerate (lines 140--144). Every definition, display and label of those lines is retained unchanged. Omitted from this package and not redistributed: 1--139, 145--173, 208--605. Nothing inside a retained excerpt is omitted or altered: every retained body is delivered exactly as the article's lines are written, so no omission receipt of any kind accompanies this package. Citations: the retained excerpts carry no citation command at all, so no bibliography entry of the article is transcribed and no reference is redistributed. Reference adaptation: none was needed and none was made; every reference inside the delivered unit resolves inside it, so no unresolved reference or citation is produced. Printed slips kept literal: no printed word, symbol, hypothesis or number of the counted statement or of its proof was altered. Layout and class: the article's own class is \texttt{amsart} with packages including amsmath, amssymb; the wrapper is the lane's pinned \texttt{amsart} editorial wrapper at 10pt on A4, which restates the theorem-like environments the retained text needs and adds the article's own macro definitions that the retained text needs (none), with extra wrapper packages (bm, bbm, dsfont, xspace, cancel, mathtools). The delivered numbering is the unit's own. Assets: no retained span contains a figure environment, a graphics-inclusion command, a PDF-page inclusion, a table or a TikZ picture; no image file is copied into this package, the package declares no asset dependency and no image file is redistributed with it. Licence: version arXiv:2602.20481v1 of this article is distributed under the Creative Commons Attribution 4.0 International licence, as recorded in the arXiv licence metadata for this exact version and shown on the abstract page https://arxiv.org/abs/2602.20481v1. A full rights notice is delivered at licenses/arxiv-2602.20481-CC-BY-4.0.txt. Characters: every retained excerpt is pure ASCII, so the delivered engine, XeLaTeX with its default Latin Modern text font, reports no missing character.
\begin{prop} \label{non-degenerate} Let $p(x)$ be an irreducible monic polynomial over $k$ with $p(0)\neq 0$ such that $p(x)\neq p^*(x)$ and $E=k[x]/(p(x)p^*(x))$. 
Suppose $(M, h)$ is a hermitian space over $E$. 

Then $(M, h)$ is non-degenerate if and only if $\det(h(e_i, e_j))_{1\leq i, j \leq n}$ is invertible in $E$ for a basis $\{ e_1, \cdots, e_n \}$ of $M$ over $E$.
\end{prop}

\begin{prop} \label{orth-base}  Let $p(x)$ be an irreducible monic polynomial over $k$ with $p(0)\neq 0$ such that $p(x)\neq p^*(x)$ and $E=k[x]/(p(x)p^*(x))$.  Assume that $char(k)\neq 2$. 

If $(M, h)$ is a non-degenerate hermitian space over $E$ with respect to the involution induced by $x\mapsto x^{-1}$, then $M$ has an orthogonal basis $\{ v_1, \cdots, v_n \}$ over $E$ such that $$ h(v_1, v_1) = \cdots = h(v_n, v_n) =1. $$ In particular,  any two non-degenerate hermitian spaces of the same rank over $E$ with respect to the involution induced by $x\mapsto x^{-1}$ are isomorphic. 
\end{prop}

\begin{proof} Let $\{ e_1, \cdots, e_n \}$ be a basis of $M$ over $E$. 

If $h(e_1, e_1)$ is invertible in $E$, then $M= E e_1 \perp N$ where 
$$ N= E(e_2- \frac{h(e_1, e_2)}{h(e_1, e_1)} e_1)+  \cdots + E(e_n- \frac{h(e_1, e_n)}{h(e_1, e_1)} e_1). $$
Since $\det (M) = \det (Ee_1) \cdot \det (N)$, one obtains that $N$ is also a non-degenerate hermitian space over $E$ by Proposition \ref{non-degenerate}. The result follows from induction.

Otherwise, one has that $h(e_1, e_1)$ is a zero divisor of $E$. Then there is $2\leq i_0\leq n$ such that $h(e_1, e_{i_0})\neq 0$ by the assumption that $\det (h(e_i, e_j))$ is invertible in $E$. Without loss of generality, we simply assume that $h(e_1, e_2)\neq 0$.  

Claim: There is $\lambda\in E$ such that $h(\lambda e_1+ e_2, \lambda e_1+  e_2)$ is invertible in $E$. 

Indeed, since 
$$ E\cong k[x]/(p(x)) \times k[x]/(p^*(x)) \cong K\times K $$
as $k$-algebras where $K=k(\alpha)$ and $\alpha$ is a root of $p(x)$, the involution of $E$ induced by $x\mapsto x^{-1}$ gives the involution of $K\times K$ by interchanging two coordinates.  Under this identification, one can write 
$$ h(e_1, e_1)= (a_1, a_2), \ \ \ h(e_1, e_2) = (c_1, c_2) \ \ \ \text{and} \ \ \ h(e_2, e_2) =(b_1, b_2)$$  
where $a_i, b_i, c_i\in K$ for $i=1, 2$.  Write $\lambda= (\lambda_1, \lambda_2)$. Then 
$$ h(\lambda e_1+ e_2, \lambda e_1+ e_2)=(a_1\lambda_1\lambda_2+c_1\lambda_1+c_2\lambda_2+b_1, a_2\lambda_1\lambda_2+ c_1\lambda_1+c_2\lambda_2+b_2).$$
Since $h(e_1, e_1)$ is a zero divisor of $E$, one can simply assume that $a_1=0$.  By $h(e_1, e_2)= (c_1, c_2)\neq (0, 0)$, the linear form $c_1\lambda_1+c_2\lambda_2$ is not zero. 

When $a_2\neq 0$, one considers the equation 
$$ c_1\lambda_1+c_2\lambda_2+b_1=1 .$$  Then the equation 
$$ a_2\lambda_1\lambda_2+ c_1\lambda_1+c_2\lambda_2+b_2=0 $$ has at most two solutions in $K$. Since $char(k)\neq 2$, the field $K$ contains more than two elements. Therefore there are $\lambda_1, \lambda_2$ in $K$ such that 
$$ a_2\lambda_1\lambda_2+ c_1\lambda_1+c_2\lambda_2+b_2 \neq 0 . $$ This implies that $h(\lambda e_1+ e_2, \lambda e_1+  e_2)$ is invertible in $E$. 


When $a_2=0$, there are $\lambda_1, \lambda_2\in K$ such that $$c_1\lambda_1+c_2\lambda_2 \in K\setminus \{-b_1, -b_2 \}$$ by $char(k)\neq 2$. Therefore the claim follows.

By repeating the beginning arguments to the basis $\{ \lambda e_1+ e_2, e_1, e_3, \cdots, e_n \}$ of $M$, one obtains an orthogonal basis $\{\eta_1, \cdots, \eta_n \}$. Write
$$ h(\eta_1, \eta_1)= (\alpha_1, \alpha_1), \cdots, h(e_n, e_n)=(\alpha_n, \alpha_n) $$ with $\alpha_1, \cdots, \alpha_n \in K$. Then 
$$ v_1= (\alpha_1, 1)^{-1} \eta_1, \cdots, v_n = (\alpha_n, 1)^{-1} \eta_n $$ is as required. 
\end{proof}

\end{document}
